% Generated draft: requires mathematical and source review.
\documentclass[11pt]{article}
\usepackage[margin=25mm]{geometry}
\usepackage{fontspec}
\setmainfont{Latin Modern Roman}
\usepackage{amsmath,amssymb,mathtools,unicode-math}
\setmathfont{Latin Modern Math}
\usepackage{enumitem,needspace,xurl,hyperref}
\setlength{\parindent}{0pt}
\setlength{\parskip}{5pt}
\setlength{\emergencystretch}{3em}
\newcommand{\N}{\mathbb{N}}
\newcommand{\Z}{\mathbb{Z}}
\newcommand{\Q}{\mathbb{Q}}
\newcommand{\R}{\mathbb{R}}
\newcommand{\C}{\mathbb{C}}
\newcommand{\F}{\mathbb{F}}
\newcommand{\A}{\mathbb{A}}
\newcommand{\PP}{\mathbb{P}}
\newcommand{\E}{\mathbb{E}}
\newcommand{\eps}{\varepsilon}
\newcommand{\dd}{\,\mathrm{d}}
\begin{document}
\section*{UnsolvedMath ID: 30004577}
\textbf{Source problem code:} catalogue entry
\subsection{Definitions and mathematical statement}
The supplied record does not define numeracy, specify a collection of cultures or social contexts, or provide a criterion for deciding when numeracy is uniform or varies. Its stated question is:
\begin{quote}
Is numeracy a single, uniform phenomenon, or does it vary across cultures and social contexts?
\end{quote}
This is a qualitative historical and cultural question. No mathematical objects, domains, comparison relation, boundary conditions, or quantifiers are specified. Consequently, the supplied evidence does not determine a precise mathematical proposition. Defining these ingredients would require additional substantive assumptions. The intended meaning of uniformity and the forms of variation under consideration remain unresolved.
\subsection{Short English statement}
The question asks whether numeracy is uniform across cultures and social contexts or varies between them. The supplied record gives no definitions or criteria that would turn this inquiry into a precise mathematical statement.
\subsection{Sources}
\begin{itemize}
\item Local UnsolvedMath registry: supplied record with permanent ID \(30004577\), title \emph{Cultural and Social Variation in Numeracy}, source problem code \texttt{catalogue entry}, in the Oberwolfach Reports Open Problems collection. No individual proposer is supplied; the recorded proposed year is 2020.
\item Supplied report attribution: \emph{History of Mathematics: A Global Cultural Approach} (2021), Oberwolfach Reports, pp.\ 2021--2092, DOI \texttt{10.4171/owr/2020/41}. Supplied source and literature-review link: \url{https://doi.org/10.4171/owr/2020/41}. The background describes Oberwolfach Reports as documenting mathematical workshops and their open-problem discussions; no precise location for this question within the report is supplied.
\item Supplied literature assessment, checked on \texttt{2026-08-22}: the extracted text is characterized as a rhetorical fragment rather than an identifiable mathematical claim, so no proposition can be verified as solved or disproved; it remains labeled open under the supplied triage rule. This is supplied dated curation, not an independently verified literature assessment. The background requests independent verification of time-sensitive status and cited claims.
\item The registry labels the statement \texttt{exact} and reports that it is internally readable and complete and was retained byte-for-byte. These are supplied extraction qualifications, not independent verification of the report text or mathematical precision.
\end{itemize}
\end{document}
